\section*{Chapter \ref{ch:b1:cross-sections} --- Cross-Sections}

\begin{solution}{exo:b1:cross-sections:1}
At 1:25\,000, \qty{1}{cm} on the map is \qty{250}{m} on the ground. The
vertical scale is \qty{1}{cm} for \qty{100}{m}, so the vertical scale is
$250/100 = 2.5$ times the horizontal one: $e = 2.5$.
\end{solution}

\begin{solution}{exo:b1:cross-sections:2}
$\gamma = \ang{90} - \ang{60} = \ang{30}$, and $\tan\beta = \tan\ang{40}
\cos\ang{30} = 0.8391\times0.8660 = 0.7267$: $\beta = \ang{36.0}$, the bed
going down toward the east end of the section. In the section of \omterm{def:b1:spherical-astronomy:horizontal}{azimuth}
\ang{000}, $\gamma = \ang{90}$ and $\beta = 0$: the section runs along the
\omterm{def:b1:geological-maps:attitude}{strike} and the bed appears horizontal.
\end{solution}

\begin{solution}{exo:b1:cross-sections:3}
$t = h\cos\delta = 120\times\cos\ang{12.16} = 120\times0.9776 =
\qty{117.3}{m}$.
\end{solution}

\begin{solution}{exo:b1:cross-sections:4}
$\tan\ang{12} = 0.2126$. With $e = 2$: $\arctan 0.4251 = \ang{23.0}$; with
$e = 3$: $\arctan 0.6377 = \ang{32.5}$; with $e = 5$: $\arctan 1.0628 =
\ang{46.7}$. The angle grows much less than $e$ itself, and never reaches
\ang{90}.
\end{solution}

\begin{solution}{exo:b1:cross-sections:5}
(a) $t = 300\sin\ang{25} = \qty{126.8}{m}$. (b) $t = 300\sin\ang{35} =
\qty{172.1}{m}$. (c) $t = 300\sin\ang{15} = \qty{77.6}{m}$. The same \omterm{def:b1:geological-maps:geological-map}{outcrop}
width can mean thicknesses differing by more than a factor of two: the
slope of the ground must always be taken into account.
\end{solution}

\begin{solution}{exo:b1:cross-sections:6}
With \omterm{def:b1:geological-maps:attitude}{dip direction} $\alpha$ and \omterm{def:b1:geological-maps:attitude}{dip} $\delta$, the theorem gives in the
north--south section ($\theta = 0$) $\tan\ang{20} = \tan\delta\cos\alpha$,
and in the east--west section ($\theta = \ang{90}$) $\tan\ang{30} =
\tan\delta\cos(\alpha - \ang{90}) = \tan\delta\sin\alpha$. Dividing,
$\tan\alpha = 0.5774/0.3640 = 1.586$, $\alpha = \ang{057.8}$; adding the
squares, $\tan\delta = \sqrt{0.3640^2 + 0.5774^2} = 0.6825$, $\delta =
\ang{34.3}$. Both \omterm{def:b1:cross-sections:apparent-dip}{apparent dips} are smaller than the true \omterm{def:b1:geological-maps:attitude}{dip}, as they must
be.
\end{solution}

\begin{solution}{exo:b1:cross-sections:7}
$100/\tan\ang{12.16} = 100/0.2155 = \qty{464.1}{m}$. On the map, the
\omterm{def:b1:cross-sections:structure-contour}{structure contours} of the east limb are indeed a little less than half a
kilometre apart, equally spaced and parallel: the limb is a plane.
\end{solution}

\begin{solution}{exo:b1:cross-sections:8}
B1: $524.02 - 488.98 = \qty{35.04}{m}$; B3: $529.98 - 494.93 =
\qty{35.05}{m}$. Above B2, the base of the shale is the top of the
limestone, \qty{180}{m} above its base: $280.16 + 180 = \qty{460.16}{m}$,
higher than the ground (\qty{380.64}{m}), which is why it is eroded there.
Between B1 and B2 the contact rises by $460.16 - 35.04 = \qty{425.12}{m}$
over \qty{2000}{m}: a slope of $0.2126$, an \omterm{def:b1:cross-sections:apparent-dip}{apparent dip} of \ang{12.0}. B1
and B3, at the same distance on either side of B2, find every contact at the
same elevation within \qty{1}{cm}: the \omterm{def:b1:geological-maps:fold}{fold} is symmetric about a vertical
axial plane at $x = \qty{3.5}{km}$, its two limbs dipping equally, and no
fault cuts the section between them.
\end{solution}

\begin{solution}{exo:b1:cross-sections:9}
\omterm{def:b1:cross-sections:throw}{Heave} $= 150/\tan\ang{50} = 150/1.1918 = \qty{125.9}{m}$. In the section of
\omterm{def:b1:spherical-astronomy:horizontal}{azimuth} \ang{150}, $\gamma = \ang{30}$ and $\tan\beta = 1.1918\times0.8660 =
1.0321$: the fault appears with the \omterm{def:b1:cross-sections:apparent-dip}{apparent dip} \ang{45.9}. The cutoffs of
a horizontal bed are \qty{150}{m} apart vertically on that trace, hence
$150/1.0321 = \qty{145.3}{m}$ apart along the section, more than the \omterm{def:b1:cross-sections:throw}{heave},
since the section crosses the fault obliquely.
\end{solution}

\begin{solution}{exo:b1:cross-sections:10}
$\cos\gamma = \tan\ang{20}/\tan\ang{40} = 0.3640/0.8391 = 0.4338$, so
$\gamma = \ang{64.3}$: the sections of \omterm{def:b1:spherical-astronomy:horizontal}{azimuths} $\alpha \pm \ang{64.3}$
(and the opposite directions, where the bed rises). The \omterm{def:b1:cross-sections:apparent-dip}{apparent dip} is
\ang{40.0}, \ang{39.6}, \ang{38.3} and \ang{36.0} for $\gamma = 0$, \ang{10},
\ang{20}, \ang{30}, but \ang{22.8}, \ang{16.0}, \ang{8.3} and $0$ for
$\gamma = \ang{60}$, \ang{70}, \ang{80}, \ang{90}. Indeed, differentiating
$\tan\beta = \tan\delta\cos\gamma$,
$\dfrac{\mathrm d\beta}{\mathrm d\gamma} =
-\dfrac{\tan\delta\sin\gamma}{1 + \tan^2\delta\cos^2\gamma}$, which is zero
at $\gamma = 0$ and equals $-\tan\delta$ at $\gamma = \ang{90}$: a section a
little off the \omterm{def:b1:geological-maps:attitude}{dip direction} shows nearly the true \omterm{def:b1:geological-maps:attitude}{dip}, while near the \omterm{def:b1:geological-maps:attitude}{strike}
every degree of \omterm{def:b1:spherical-astronomy:horizontal}{azimuth} changes the \omterm{def:b1:cross-sections:apparent-dip}{apparent dip} by almost a degree.
\end{solution}

\begin{solution}{exo:b1:cross-sections:11}
The beds \omterm{def:b1:geological-maps:attitude}{dip} east, so the normal to the bedding, followed downward, is
tilted \ang{12.16} from the vertical toward the west, toward \ang{260.7}.
(a) A vertical borehole makes the angle $\psi = \ang{12.16}$ with it:
$176.0/\cos\ang{12.16} = \qty{180.0}{m}$, the \omterm{def:b1:cross-sections:thickness}{vertical thickness}. (b) A
borehole plunging \ang{60} toward \ang{260.7} is tilted \ang{30} from the
vertical on the same side as the normal: $\psi = \ang{30} - \ang{12.16} =
\ang{17.84}$, $\cos\psi = 0.9519$, and the core is $176.0/0.9519 =
\qty{184.9}{m}$ long. (c) Toward \ang{080.7} it is tilted on the other side:
$\psi = \ang{30} + \ang{12.16} = \ang{42.16}$, $\cos\psi = 0.7413$, and the
core is $176.0/0.7413 = \qty{237.4}{m}$ long. The core is shortest, equal to
$t$, when the borehole runs along the normal, plunging \ang{77.84} toward
\ang{260.7}.
\end{solution}

\begin{solution}{exo:b1:cross-sections:12}
In the east--west vertical plane, with the angles measured from the east,
the upper beds rise toward the east at \ang{20} and the lower beds go down
toward the east at \ang{70}. Rotating the \omterm{def:b1:geological-maps:geological-map}{outcrop} about the \omterm{def:b1:geological-maps:attitude}{strike} by
\ang{20}, east side down, makes the upper beds horizontal and brings the
lower beds to $\ang{70} + \ang{20} = \ang{90}$: they were vertical when the
upper beds were deposited. History, from the oldest event: the lower beds are
deposited flat; they are turned up on edge, to the vertical, and raised;
their edges are eroded to a surface; the upper beds are deposited across that
surface; the whole is tilted by \ang{20} down to the west; erosion cuts
today's ground surface. It is the history of Hutton's \omterm{def:b1:geological-maps:geological-map}{outcrop} at Siccar
Point.
\end{solution}

\begin{solution}{pb:b1:cross-sections:1}
\textbf{1.} The thickness above the tunnel is $z - 200$: largest at
$x = \qty{1}{km}$, $556 - 200 = \qty{356}{m}$; smallest at $x =
\qty{9}{km}$, $339 - 200 = \qty{139}{m}$.

\textbf{2.} \qty{1}{km} of elevation measures \qty{4.2}{cm} and \qty{1}{km}
of distance \qty{1.4}{cm}: $e = 4.2/1.4 = 3$.

\textbf{3.} $380.64 - 220.48 = \qty{160.16}{m}$.

\textbf{4.} In B2 the base of the limestone is at $380.64 - 100.48 =
\qty{280.16}{m}$ and the base of the sandstone at \qty{160.16}{m}: the
tunnel, at \qty{200}{m}, is in the sandstone.

\textbf{5.} B1: $524.02 - 668.98 = \qty{-144.96}{m}$; B3: $529.98 - 674.93 =
\qty{-144.95}{m}$.

\textbf{6.} From west to east: chalk, shale, limestone, sandstone (in the core
of the \omterm{def:b1:geological-maps:fold}{anticline}), limestone, shale, and chalk again as far as
$x = \qty{9}{km}$, through the \omterm{def:b1:geological-maps:fold}{syncline} and across the fault.

\textbf{7.} $\gamma = \ang{80.7} - \ang{90} = \ang{-9.3}$; $\tan\beta =
\tan\ang{12.16}\cos\ang{9.3} = 0.2155\times0.9869 = 0.2126$: $\beta =
\ang{12.0}$, down toward the east.

\textbf{8.} $(280.16 + 144.95)/2000 = 425.11/2000 = 0.2126$: \ang{12.0}.

\textbf{9.} $\arctan(3\times0.2126) = \arctan 0.6377 = \ang{32.5}$.

\textbf{10.} $\gamma = \ang{120} - \ang{90} = \ang{30}$; $\tan\beta =
\tan\ang{60}\cos\ang{30} = 1.732\times0.866 = 1.500$: $\beta = \ang{56.3}$.

\textbf{11.} The limbs \omterm{def:b1:geological-maps:attitude}{dip} toward \ang{080.7} or \ang{279.3}, only
\ang{9.3} from the line of section, and $\cos\ang{9.3} = 0.987$; the fault
\omterm{def:b1:geological-maps:attitude}{dips} toward \ang{120}, \ang{30} from it, and $\cos\ang{30} = 0.866$.

\textbf{12.} The fault plane is $z = 250 - \tan\ang{60}\,[(x - 4300)
\sin\ang{120} + (y - 1000)\cos\ang{120}]$. With $y = 6000$, $(y - 1000)
\cos\ang{120} = -2500$, so $(x - 4300)\sin\ang{120} = 2500 - (z - 250)/
\tan\ang{60}$, and $1/\tan\ang{60} = \tan\ang{30}$.

\textbf{13.} At $z = 200$: $x = 4300 + (2500 + 50\times0.5774)/0.8660 =
4300 + 2528.9/0.8660 = \qty{7220}{m}$, that is \qty{7.22}{km} from C.

\textbf{14.} \omterm{def:b1:cross-sections:throw}{Heave} $= 200/\tan\ang{60} = \qty{115.5}{m}$; separation along
C--D $= 200/1.500 = \qty{133.3}{m}$.

\textbf{15.} The fault \omterm{def:b1:geological-maps:attitude}{dips} east-south-east and is normal: the hanging wall,
to the east, is downthrown. At tunnel level the chalk lies on both sides of
the fault: west of it, the chalk of the \omterm{def:b1:geological-maps:fold}{syncline}; east of it, chalk brought
down \qty{200}{m}.

\textbf{16.} The tunnel stays in chalk across the fault, so no change of rock
at the face announces it; only a zone of crushed rock, and perhaps water,
does. The engineers predict its position from the section (here
\qty{7.22}{km}), drill probe holes ahead of the machine as it approaches, and
prepare support for the crushed zone.

\textbf{17.} $668.98 - 488.98 = \qty{180.00}{m}$.

\textbf{18.} $t = 180\cos\ang{12.16} = \qty{176.0}{m}$.

\textbf{19.} On the crest ($x = 3500$) the base of the sandstone is at
\qty{160.16}{m} (question 3); it goes down on both limbs with the \omterm{def:b1:cross-sections:apparent-dip}{apparent
dip} \ang{12.0} of question 7, that is by $\tan\ang{12.0}$ per metre of
horizontal distance from the axis.

\textbf{20.} The base of the limestone is \qty{120}{m} above the base of the
sandstone, at $280.16 - u\tan\ang{12.0}$ with $u = |x - 3500|$, and its top
\qty{180}{m} higher, at $460.16 - u\tan\ang{12.0}$. With $\tan\ang{12.0} =
0.21256$: the base is below \qty{200}{m} when $u > 80.16/0.21256 = 377.1$,
and the top above it when $u < 260.16/0.21256 = 1224.0$.

\textbf{21.} $x = 3500 \pm 377.1$ and $3500 \pm 1224.0$: the tunnel is in
the limestone from \qty{2276}{m} to \qty{3123}{m} and from \qty{3877}{m} to
\qty{4724}{m}.

\textbf{22.} Along a horizontal line, the base and the top of the limestone,
\qty{180}{m} apart vertically and both sloping at $\tan\beta$, are crossed a
horizontal distance $h/\tan\beta$ apart: $180/0.21256 = \qty{846.8}{m}$ on
each limb, the same to rounding as $1224.0 - 377.1 = 846.9$.

\textbf{23.} $2\times846.8 = \qty{1693.7}{m}$: the tunnel bores through
\textbf{about \qty[text-series-to-math=true]{1694}{m} of limestone}, more than a fifth of its
\qty{8}{km} length, in two stretches of \qty{847}{m} on either side of the
sandstone core of the \omterm{def:b1:geological-maps:fold}{anticline}.
\end{solution}
